\documentclass[journal]{IEEEtran}

\usepackage[T1]{fontenc}
\usepackage{amsmath,amsfonts}
\usepackage{newtxtext}
\usepackage{newtxmath}

\usepackage{algorithmic}

\usepackage{cite}
\usepackage{graphicx}
\usepackage{booktabs}

\usepackage[hidelinks]{hyperref}

\begin{document}

\title{title}

\author{Andy Babcock \\
    \textit{Steve Sanghi College of Engineering, Northern Arizona University} \\
    Flagstaff, AZ, USA \\
    aab726@nau.edu
}

\maketitle
Tests in school and college are high-stress environments. Even when properly prepared, it is common to feel nervous while taking a test, which can potentially reduce performance. Then, when someone with higher than average anxiety levels, it can become difficult to accurately demonstrate that they actually have learned the content of the course. Because test-taking requires safeguards to prevent cheating, like working in a classroom with nothing aside from a pencil or calculator, there are limited ways for students to manage their anxiety. Various kinds of noise, like music or white noise could mitigate anxiety and in turn improve test scores. From existing works that cover various kinds of noise in high-stress environments, we would initially expect white noise to make no significant difference in stress level or performance, while music may lower stress and increase performance~\cite{ShahrokhiShabnam2022Teow}.

What is the effect of noise on test-taking anxiety and performance?

Listening to white noise shows no statistically significant reduction in anxiety or improvement in test scores, tested with a wearable EEG measuring anxiety. If there is a reduction in anxiety, there may also be an improvement in test scores. Additionally, we will test familiar and unfamiliar music, with the expectation that familair music reduces anxiety and improves test scores, while unfamiliar music reduces anxiety at the cost of lower test scores. Alternatively, familiar music may not improve test scores, or unfamilir music may improve test scores similarly to familiar music. Additionally, unfamiliar music can be split into distinct genres that may have varying levels of distraction: jazz, for example, may be less distracting than pop or rock music.

Using a wearable EEG that measures stress, we will test these claims by repeatedly taking tests while various kinds of noise play. EEGs have already proven to be 

\bibliographystyle{IEEEtran}
\bibliography{references}

\end{document}
